\documentclass[11pt,a4paper]{article}
\usepackage[margin=2.4cm]{geometry}
\usepackage{amsmath,amssymb,amsthm,booktabs,graphicx,hyperref,longtable}
\usepackage[T1]{fontenc}
\hypersetup{colorlinks=true,linkcolor=blue,citecolor=blue,urlcolor=blue}
\newtheorem{proposition}{Proposition}
\newtheorem{remark}{Remark}
\newtheorem{hypothesis}{Hypothesis}
\newtheorem{definition}{Definition}

\title{Load-Dependent Actuator Bandwidth Limits Calibration Transfer\\
in Two-Axis Thrust-Vector-Controlled Rocket Analogues}
\author{Robert Bulau\\
\small Colegiul Na\c{t}ional ``Mihai Viteazul'', Ploie\c{s}ti, Romania\\
\small \texttt{robertbulau8@gmail.com}}
\date{Revision 2 --- \today}

\begin{document}
\maketitle

\begin{abstract}
Small electric thrust-vector-control (TVC) analogues are usually designed with servo models
identified on an unloaded bench. We ask whether that calibration transfers: does an unloaded
first-order servo fit predict the closed-loop stability boundary of a reaction-loaded, reconfigured
two-axis gimbal? We first tighten the analytical baseline. A corrected planar rigid-body reduction
retains the direct lateral force produced by gimbal deflection, giving a right-half-plane transmission
zero at exactly the normalized frequency $1$. We prove that the closed-loop spectrum of the declared
PD cascade depends on exactly five dimensionless groups, in which the inertia ratio $\chi$ never
appears alone; consequently ``geometry sets the boundary'' is a statement about gain convention, not
physics. We then derive a closed-form ceiling on the normalized position gain in the ideal-actuator
limit, showing that the commonly quoted necessary condition $\chi k_x<1$ is loose by roughly an order
of magnitude and that the true ceiling is the non-minimum-phase crossover limit
$\omega_{\text{out}}\lesssim 0.3\,\omega_0$. Two distinct instability mechanisms bound the stable
region: an actuator-bandwidth boundary and this gain ceiling. The experimental contribution is a
single falsifiable claim. We hypothesise that reflected load inertia is the dominant, and largest,
source of calibration-transfer error in hobby-class digital servos, and that a one-parameter
correction $\omega_a(J_{\text{refl}})$ recovers held-out boundary prediction. The protocol is a
reaction-loaded two-axis gimbal with independent magnetic-encoder truth, swept over calibrated
inertia discs, closed in hardware-in-the-loop around the derived cascade, with complete
configurations held out. Deliverables are an open dataset, a reproducible pipeline, and a quantified
transfer envelope. All numerical values reported here are synthetic implementation checks; no
hardware measurement is claimed.
\end{abstract}

\noindent\textbf{Keywords:} thrust vector control; actuator identification; load-dependent bandwidth;
cascaded control; model transfer; hardware-in-the-loop.

\section{Introduction}

Electric rocket analogues make thrust vectoring repeatable without combustion modelling, and the
literature on stabilising them is mature \cite{spannagl2021,linsen2022,erocket2026,quadrocket2026}.
The question that remains open at this scale is narrower and more practical. Almost every small TVC
build identifies its servos on a bench with nothing attached, then uses that model to design a
controller for a gimbal that carries a motor, a propeller, and a mounting bracket. Nobody has
measured how much that substitution costs.

At large scale the effect is known to exist. Orr et al.\ \cite{orr2023} document actuator--structure
interaction on SLS, where loaded actuator response departs materially from isolated calibration. No
comparable measurement exists for the hobby-class digital servos that every small TVC testbed uses,
and the relevant physics is not obviously the same: an SLS electrohydraulic actuator and a
\$20 bus servo with a plastic-and-metal gear train fail for different reasons.

This paper makes two contributions. Section~\ref{sec:analysis} tightens the analytical baseline for
a declared PD cascade, and in doing so extracts three results that the prior small-scale TVC
literature states loosely or not at all: an exact five-group reduction, a closed-form
non-minimum-phase gain ceiling, and the two-mechanism structure of the stable region.
Section~\ref{sec:experiment} states and operationalises one falsifiable experimental hypothesis about
why unloaded calibration fails, together with the protocol that can reject it.

\subsection{What is explicitly not claimed}

Newton--Euler modelling, nondimensionalisation, PD cascades, Hurwitz tests, dynamic similarity,
servo identification, rate limiting, and gimbal Jacobians are all standard. Gain-region stability
maps date to at least 1969 \cite{jacqmein1969}. First-order servo lag with rate limiting is already
explicit in small-launcher simulation \cite{santos2024}. Two-axis gimbal construction and response
characterisation are covered \cite{singh2025,erocket2026}, as is actuator-level nonlinear
compensation on dual-actuator platforms \cite{pidnet2026,ddobs2026}. This paper does not claim any of
them. It claims a measured parametric law for how one specific, universally used calibration
shortcut degrades, and the resulting transfer envelope.

\section{Prior art and the remaining gap}

Table~\ref{tab:prior} summarises the closest work. The audit was conducted over arXiv, NASA/NTRS,
institutional repositories, publisher pages, and indexed patent text. It is a public-source audit,
not an authenticated exhaustive search, and absence of an exact title is never treated as evidence
of originality.

\begin{longtable}{p{3.0cm}p{6.0cm}p{6.0cm}}
\caption{Closest prior work and the remaining distinction.}\label{tab:prior}\\
\toprule
\textbf{Work} & \textbf{System and method} & \textbf{Consequence here} \\
\midrule
\endfirsthead
\toprule
\textbf{Work} & \textbf{System and method} & \textbf{Consequence here} \\
\midrule
\endhead
Jacqmein et al.\ 1969 \cite{jacqmein1969} &
Hypothetical lunar flying vehicle; Hurwitz-minor stability maps in vehicle parameters. &
Parameter stability maps are old. Only the measured load-transfer law is new. \\
Hauser et al.\ 1992 \cite{hauser1992} &
V/STOL force--moment coupling, slightly non-minimum-phase dynamics. &
The RHP zero is established physics. Prop.~\ref{prop:ceiling} quantifies its gain cost in this
normalisation; discovery is not claimed. \\
Corr\^{e}a 2013 \cite{correa2013} &
Single-axis PD gain regions via D-decomposition under gimbal limits. &
Gain-region tuning alone is too close. The distinction is held-out, load-varied prediction. \\
Spannagl et al.\ 2021 \cite{spannagl2021} &
Low-cost electric rocket, MPC guidance, flight demonstration. &
Construction and GNC are infrastructure, not novelty. \\
Linsen et al.\ 2022 \cite{linsen2022} &
Two-axis gimbal, contra-rotating props, NMPC, EKF actuator-offset estimation. &
Estimating actuator offset online is not the same as predicting a boundary offline from calibration. \\
Orr et al.\ 2023 \cite{orr2023} &
SLS engine--actuator--structure interaction, validated against full-scale test. &
Establishes that loading matters at large scale. Motivates, but does not answer, the small-scale
question. \\
Santos \& Oliveira 2024 \cite{santos2024} &
Six-DOF launcher simulation with first-order servos and rate limits. &
Actuator lag is already modelled in the literature; we must not claim otherwise. Its $\omega_a$
is assumed constant, which is the assumption tested here. \\
Singh 2025 \cite{singh2025} &
Low-cost printed two-axis gimbal, response assessed from video. &
Construction is covered. Video-derived response duration is not an identified time constant under load. \\
Fonte et al.\ 2025 \cite{fonte2025} &
Simulated 2-D electric rocket, ideal actuator. &
An idealised study motivating exactly this measurement. \\
QuadRocket 2026 \cite{quadrocket2026} &
Quadrotor-as-TVC actuator, adaptive backstepping, actuator dynamic-surface controller. &
Actuation dynamics are addressed, but the mechanism avoids servo-gimbal gear-train loading. \\
PIDNet 2026 \cite{pidnet2026} &
Dual linear actuator gimbal, neural compensation, two-axis experiments. &
A compensator, not a calibration-transfer benchmark. \\
DD-DOB 2026 \cite{ddobs2026} &
Data-driven input-matrix adaptation, dual-channel actuator HIL. &
Adapts online to unknown actuation; does not quantify how far offline calibration carries. \\
\bottomrule
\end{longtable}

Patent search over TVC, gimbal stability, bandwidth, and adaptive effectiveness returned
EP3529683B1 \cite{ep3529683} (two-axis thrust-vectoring multicopter hardware with model-based
decoupling) and WO2025122161A1 \cite{wo2025122161} (operating-point-based gain adjustment for
suspended loads). Both foreclose broad claims to two-axis gimbal control or load-adaptive gains.
Neither states a measured calibration-transfer envelope. This is a search observation, not legal
clearance.

\textbf{The gap.} Every inspected source either assumes a constant actuator model, identifies it
online, or compensates for it adaptively. None measures the offline question: given a calibration at
one load, how far does the predicted stability boundary remain accurate as load changes, and what
minimal correction restores it. That question is small, but it is answerable in six weeks on a bench,
and its answer is immediately usable by everyone building these platforms.

\section{Corrected dynamics and analytical baseline}\label{sec:analysis}

\subsection{Rigid body}

Let inertial $e_3$ point upward, $R\in SO(3)$ be the body-to-inertial rotation, and $p,v,\omega$ be
centre-of-mass position, velocity, and body angular rate. With gimbal angles $(\alpha,\beta)$,
\begin{equation}
R_g(\alpha,\beta)=R_x(\alpha)R_y(\beta),\qquad
f_B = T R_g e_3 = T\begin{bmatrix}\sin\beta\\ -\sin\alpha\cos\beta\\ \cos\alpha\cos\beta\end{bmatrix}.
\end{equation}
For thrust application point $r=-\ell e_3$, $\ell>0$,
\begin{equation}
\dot p = v,\quad m\dot v = -mge_3 + Rf_B + F_{\text{ext}},\quad
\dot R = R[\omega]_\times,\quad J\dot\omega + \omega\times J\omega = r\times f_B + \tau_{\text{prop}} + \tau_{\text{ext}}.
\end{equation}
The force is rotated by $R_g$ once; the moment uses the same body-frame force. Propeller reaction and
gyroscopic terms are retained symbolically and bounded in Sec.~\ref{sec:gyro} rather than silently
dropped. The longitudinal axis is called roll; the two controlled channels are transverse tilt,
$q=[\theta_y,-\theta_x]^{\mathsf T}$ at small angles. Software axis names must follow these
coordinates, not the conventions of any cited paper.

\subsection{Planar reduction}

For rotation $R_y(\theta)$ with physical deflection $\delta=-u$, the thrust angle from vertical is
$\theta-u$ and the moment is $+T\ell\sin u$. With $a=T/m$, $b=T\ell/J$,
\begin{equation}
\ddot x = a\sin(\theta-u),\qquad \ddot\theta = b\sin u,\qquad \dot u = \omega_a(u_c-u),
\end{equation}
linearising to $\ddot x = a(\theta-u)$, $\ddot\theta = bu$. The term $-au$ is the direct lateral force
required to generate the corrective moment. Retaining it,
\begin{equation}
\frac{x(s)}{u(s)} = \frac{a(b-s^2)}{s^4},
\end{equation}
with a right-half-plane zero at $+\sqrt b$. Dropping the direct force removes it and makes every
subsequent gain bound wrong in the optimistic direction.

With $\omega_0=\sqrt b$, $\tau=\omega_0 t$, $X=x/\ell$, $\chi=J/(m\ell^2)$, $\eta=\omega_a/\omega_0$,
and primes denoting $d/d\tau$,
\begin{equation}\label{eq:dimless}
X'' = \chi(\theta-u),\qquad \theta''=u,\qquad u'=\eta(u_c-u).
\end{equation}
In these coordinates the RHP zero sits at exactly $s=+1$: the non-minimum-phase obstruction is
located at $\omega_0=\sqrt{T\ell/J}$ for every configuration, which is what makes it a usable design
scale. Note that $\omega_0$ is an authority scale, not a natural frequency of the uncontrolled body.
At upright hover $T=mg$, so a thrust sweep must vary mass consistently, use a non-hover trim, or
model the rig.

\subsection{Controller and the five-group reduction}

Define the normalised PD cascade
\begin{equation}\label{eq:cascade}
\theta_c = -k_x X - k_v X',\qquad u_c = k_\theta(\theta_c-\theta) - k_q\theta',
\end{equation}
equivalently $\theta_c = -(k_x/\ell)x - (k_v/\ell\omega_0)\dot x$ and
$u_c = k_\theta(\theta_c-\theta)-(k_q/\omega_0)\dot\theta$ in physical variables. With state
$\xi=[X,X',\theta,\theta',u]^{\mathsf T}$,
\begin{equation}\label{eq:A}
A=\begin{bmatrix}
0&1&0&0&0\\ 0&0&\chi&0&-\chi\\ 0&0&0&1&0\\ 0&0&0&0&1\\
-\eta k_\theta k_x & -\eta k_\theta k_v & -\eta k_\theta & -\eta k_q & -\eta
\end{bmatrix},
\end{equation}
\begin{equation}\label{eq:poly}
P(s)=s^5+\eta s^4+\eta(k_q-\chi k_\theta k_v)s^3+\eta k_\theta(1-\chi k_x)s^2+\eta\chi k_\theta k_v s+\eta\chi k_\theta k_x.
\end{equation}

\begin{proposition}[Five-group reduction]\label{prop:five}
The closed-loop spectrum of \eqref{eq:A} is a function of exactly five dimensionless groups,
\[
\Pi=\{\eta,\;k_\theta,\;k_q,\;g_x:=\chi k_x,\;g_v:=\chi k_v\},
\]
and the inertia ratio $\chi$ never enters independently of the outer-loop gains.
\end{proposition}
\begin{proof}
Substituting $k_x=g_x/\chi$, $k_v=g_v/\chi$ into \eqref{eq:poly} eliminates $\chi$ identically,
giving $P(s)=s^5+\eta s^4+\eta(k_q-g_v k_\theta)s^3+\eta k_\theta(1-g_x)s^2+\eta g_v k_\theta s+\eta g_x k_\theta$.
Equivalently, the change of variable $Z=X/\chi$ in \eqref{eq:dimless} yields $Z''=\theta-u$ with the
outer loop reading $\theta_c=-g_xZ-g_vZ'$.
\end{proof}

\begin{remark}
Proposition~\ref{prop:five} has a direct experimental consequence that is easy to get wrong. Holding
\emph{normalised gains} $k_x,k_v$ fixed across configurations is not the same as holding the
\emph{dimensionless closed loop} fixed: changing inertia moves $\chi$, hence $g_x$ and $g_v$, hence
the spectrum. A cross-configuration study is therefore interpolation and extrapolation in a
five-dimensional $\Pi$-space, not a similarity collapse. Any trajectory-overlay figure that fixes
$\chi$ demonstrates the identity of two identical ODEs and constitutes an implementation check only.
\end{remark}

\subsection{Two instability mechanisms}

\begin{proposition}[Attitude subsystem]\label{prop:att}
For $\eta,k_\theta,k_q>0$, the attitude--actuator subsystem with zero attitude reference is
asymptotically stable if and only if $\eta k_q>k_\theta$.
\end{proposition}
\begin{proof}
Its polynomial is $s^3+\eta s^2+\eta k_q s+\eta k_\theta$; the cubic Routh condition
$\eta(\eta k_q)>\eta k_\theta$ gives the result.
\end{proof}

\begin{proposition}[Non-minimum-phase gain ceiling]\label{prop:ceiling}
In the ideal-actuator limit $\eta\to\infty$, with $a_1=k_q-g_vk_\theta$ and $a_3=g_vk_\theta$, the
closed loop is asymptotically stable only if
\begin{equation}\label{eq:ceiling}
g_x \;<\; g_x^{\max} \;=\; \frac{a_3(a_1k_\theta-a_3)}{k_\theta\,a_1\,k_q}.
\end{equation}
Equivalently, the nominal outer-loop frequency obeys
$\omega_{\text{out}}/\omega_0=\sqrt{g_x}<\sqrt{g_x^{\max}}$.
\end{proposition}
\begin{proof}
As $\eta\to\infty$, dividing \eqref{eq:poly} by $\eta$ and discarding $s^5/\eta$ gives the quartic
$s^4+a_1s^3+a_2s^2+a_3s+a_4$ with $a_2=k_\theta(1-g_x)$, $a_4=g_xk_\theta$. The quartic Routh
condition $a_1a_2a_3>a_3^2+a_1^2a_4$ rearranges to \eqref{eq:ceiling} using $a_1+a_3=k_q$.
\end{proof}

For $k_\theta=1$, $k_q=2$, $g_v=0.2$ this gives $g_x^{\max}=0.0\overline{8}$, so
$\omega_{\text{out}}<0.298\,\omega_0$. The necessary condition $g_x<1$ obtained from coefficient
positivity is therefore loose by a factor of $11$, and $\omega_{\text{out}}<0.3\,\omega_0$ is the
correct rule of thumb --- recognisably the standard ``crossover below a third of the RHP zero''
limit, here derived exactly for this cascade. Two mechanisms consequently bound the stable region:
the actuator-bandwidth boundary of Prop.~\ref{prop:att}, moved by outer-loop feedback, and the
$\eta$-independent ceiling \eqref{eq:ceiling}. Table~\ref{tab:mech} shows the crossover.

\subsection{Two-axis model}

With $u=Gz$ the local corrective deflection ($u=[-\beta,\alpha]^{\mathsf T}$ at ideal trim),
$\ddot x = a(q-Gz)$, $\ddot q = BGz$, $B=\mathrm{diag}(T\ell/J_y,\,T\ell/J_x)$. Taking
$\omega_*^2=T\ell/J_y$, $\gamma=\chi_y$, $\bar B=B/\omega_*^2$, $E=\Lambda/\omega_*$, and nominal
inverse mixing $\widehat G^{-1}$,
\begin{equation}
q_c=-K_xX-K_vV,\qquad z_c=\widehat G^{-1}\{K_\theta(q_c-q)-K_qQ\},
\end{equation}
\begin{equation}
A=\begin{bmatrix}
0&I&0&0&0\\ 0&0&\gamma I&0&-\gamma G\\ 0&0&0&I&0\\ 0&0&0&0&\bar BG\\
-L_x&-L_v&-L_\theta&-L_q&-E\end{bmatrix},\quad
L_\bullet = E\widehat G^{-1}K_\theta K_\bullet .
\end{equation}
A scalar $\mathrm{cond}(G)$ cannot summarise this: the identity and a $60^\circ$ rotation share
condition number one yet present different feedback directions to an uncompensated controller. Note,
however, that with exact static mixing $\widehat G=G$ and an equal-axis plant, the rotation leaves
the local spectrum invariant; the counterexample therefore demonstrates what a condition number fails
to encode, not that misalignment is always destabilising.

\subsection{Gyroscopic coupling: a predicted, not fished-for, cross-axis effect}\label{sec:gyro}

Two transverse axes do not guarantee coupling; an ideal symmetric mechanism decouples at small
angles. This study therefore does not search for coupling, it predicts one term and tests the
prediction. A propeller of polar inertia $J_p$ spinning at $\Omega$ carries angular momentum
$h=J_p\Omega e_3^{(b)}$, contributing $-\omega\times h$ to \eqref{eq:A} and coupling pitch rate into
yaw torque with normalised strength
\begin{equation}
\sigma = \frac{J_p\Omega}{J\,\omega_*}.
\end{equation}
$J_p$ and $\Omega$ are directly measurable, so $\sigma$ is a prior prediction with an uncertainty
interval. For a contra-rotating pair the net $h$ cancels and $\sigma\approx 0$ --- itself a testable
claim. This converts the two-axis hypothesis from ``measure whether coupling exists'' into a
falsifiable quantitative check.

\subsection{Finite amplitude and discrete time}

Rate limits, angle bounds, deadband, sampling, and filters do not compress into \eqref{eq:poly}. A
hard rate limit typically leaves the infinitesimal linearisation unchanged while dominating
finite-amplitude behaviour. Additional groups are $r_i=\dot z_{\max,i}/(\omega_*z_*)$,
$d_i=z_{\text{dead},i}/z_*$, $h=\omega_*T_s$, $d_t=\omega_*t_d$, $f_i=\omega_{f,i}/\omega_*$.
This revision makes discrete time primary: the plant is zero-order held and the closed loop
formed with the actual sampled controller, so $h$ is a modelled parameter rather than a caveat, and
near-boundary discretisation cannot be confused with the effect under test. Map labels have two
layers: analytical (Hurwitz / non-Hurwitz / near-marginal, with a declared tolerance) and operational
(pass / fail / ambiguous against declared finite-duration thresholds). A damped oscillation is
asymptotically stable; an abort is a censored observation, not proof of divergence.

\section{Experimental study}\label{sec:experiment}

\subsection{The claim}

\begin{hypothesis}[Primary, H1]\label{h1}
Reflected load inertia is the dominant source of calibration-transfer error for hobby-class digital
servos in this application. Specifically, a first-order model identified unloaded
($\omega_a^{(0)}$) systematically over-predicts loaded bandwidth, and the identified bandwidth
follows a one-parameter law
\begin{equation}\label{eq:law}
\omega_a(J_{\text{refl}}) = \frac{\omega_a^{(0)}}{1+J_{\text{refl}}/J_c},
\end{equation}
with a single fitted $J_c$ per servo, to within repeatability noise across the tested load range.
\end{hypothesis}

\begin{hypothesis}[Transfer, H2]\label{h2}
Substituting \eqref{eq:law} into the cascade reduces held-out stability-boundary prediction error
relative to the unloaded model by a margin exceeding measurement uncertainty.
\end{hypothesis}

\begin{hypothesis}[Cross-axis, H3]\label{h3}
Measured cross-axis dynamics improve joint-excitation prediction relative to an independent-axis
model, by an amount consistent with the prior prediction $\sigma$ of Sec.~\ref{sec:gyro}. H3 may be
rejected without invalidating H1.
\end{hypothesis}

Equation~\eqref{eq:law} is the simplest physically motivated form (a fixed-torque actuator driving
increased reflected inertia through a fixed gear ratio). It is a hypothesis to be rejected, not an
assumption. If a delay or second-order term dominates instead, that is a reportable result; the
protocol below identifies which, because candidate enrichments are selected on held-out
\emph{calibration} response and never on final stability labels.

\subsection{Why reaction loading is mandatory}

An unloaded gimbal driving a virtual plant can only fail through servo nonlinearity, loop latency,
and static mixing error --- all three directly measurable during calibration, so a richer model would
win by construction and H2 would be a tautology. The physics under test is removed by the very
fixture that makes the experiment convenient. This revision therefore makes reaction loading a
requirement rather than a preference, with calibration repeated under each tested load. This is the
single most important change from the previous draft.

\subsection{Apparatus}

A two-axis gimbal carries a brushless motor and propeller on the inner ring. Each axis is driven by a
digital bus servo and independently instrumented with a 14-bit magnetic encoder reading true joint
angle, so commanded angle is never used as truth. Reflected inertia is varied by a set of machined
discs of known moment of inertia mounted on the inner ring; a reaction torque cell records the load
actually seen by the gear train. Control runs on a microcontroller with a deterministic loop at
$1/T_s$ with logged timestamps, so jitter is measured rather than assumed. Body states are simulated
(hardware-in-the-loop): real measured gimbal motion drives a virtual free-body plant, and the
controller closes around simulated body states. This validates the influence of real loaded actuator
response on the cascade, \emph{not} physical rocket translation --- a limitation stated wherever
results appear.

\subsection{Protocol}

\begin{enumerate}
\item \textbf{Actuator identification.} For each of $\geq 6$ calibrated inertia loads per axis, excite
with logarithmic chirps and steps at $\geq 2$ amplitudes in both directions, in randomised order.
Fit $\omega_a$, and competing enrichments (pure delay, second order), by frequency-response and
prediction-error methods. Report residuals on held-out \emph{calibration} records.
\item \textbf{Law fitting.} Fit $J_c$ in \eqref{eq:law} on a training subset of loads. Hold out at
least two complete loads.
\item \textbf{Boundary estimation.} Close the cascade in HIL. Scan a software gain --- $k_x$ or $k_q$,
never $\eta$, since servo bandwidth is not continuously tunable and command filtering is an ablation,
not a different servo. At each grid point, estimate the dominant closed-loop pole pair from free
response by log-envelope and Prony/ERA fits, and regress the estimated spectral abscissa against the
scan gain. Interpolate the zero crossing and bootstrap a confidence interval over repeats. The
ambiguous band is defined as the region where the abscissa CI straddles zero.
\item \textbf{Held-out prediction.} For each held-out load, predict the boundary from the unloaded
model and from \eqref{eq:law}, then measure it. Run pitch-only, yaw-only, and joint excitation.
\item \textbf{Ablations.} Ideal-actuator model; raw-dimensional-gain transfer instead of normalised;
contra-rotating versus single propeller for $\sigma$.
\end{enumerate}

Predicting the \emph{abscissa} over the whole grid, rather than a binary stable/unstable label at a
boundary, is deliberate: it uses every run, it degrades gracefully near marginality, and it yields
per-point uncertainty. Boundary-coordinate error is then a derived secondary outcome.

\subsection{Outcomes and thresholds}

Primary: held-out boundary-coordinate error and abscissa prediction RMSE, with CI coverage and
false-stable rate outside the ambiguity band. Secondary: damping/growth rate, tracking RMSE,
saturation fraction, cross-axis residual, latency and jitter, actuator prediction error. Preregistered
planning targets are boundary error below $10\%$, false-stable predictions below $5\%$ outside the
ambiguity band, and an H2 improvement exceeding measurement uncertainty. These are proposed targets
to check against attainable precision, not observed performance; outcomes are reported whether or not
they are met.

The study is designed so that both outcomes are publishable. If \eqref{eq:law} holds, the deliverable
is a calibration correction that anyone building a small TVC platform can apply in an afternoon. If
it fails, the deliverable is the identification of which enrichment does govern transfer at this
scale, plus the dataset that shows it. What is \emph{not} publishable, and must not be dressed up, is
a parameter sweep of an ideal model.

\subsection{Statistical power}

The unit of analysis is the configuration; excitation types and amplitudes are nested within it.
Three configurations cannot support an envelope claim. The design targets $\geq 6$ load
configurations per axis spanning a factor of $\geq 4$ in reflected inertia, analysed with
configuration as a random effect. With fewer, results are reported as case studies with per-case
uncertainty and the word ``envelope'' is not used.

\section{Reproducible numerical verification}

All values below are synthetic checks of the derivations, not measurements.
For $\chi=1$, $k_x=0.02$, $k_v=0.2$, $k_\theta=1$, $k_q=2$, the full-cascade boundary is
$\eta^\star=0.602676$, against the attitude-only bound $\eta>k_\theta/k_q=0.5$: outer-loop feedback
raises the actuator-bandwidth requirement. At $\eta=0.5$ the largest real part is $+0.044274$; at
$\eta=1$ it is $-0.131262$.

\begin{table}[h]\centering
\caption{Two mechanisms bounding the stable region ($\chi=1$, $k_v=0.2$, $k_\theta=1$, $k_q=2$).}
\label{tab:mech}
\begin{tabular}{lll}
\toprule
$g_x=\chi k_x$ & $\eta^\star$ & Active mechanism \\
\midrule
0.005  & 0.61950 & actuator bandwidth ($\Delta_3$) \\
0.020  & 0.60268 & actuator bandwidth ($\Delta_3$) \\
0.050  & 0.56720 & actuator bandwidth ($\Delta_3$) \\
0.080  & 0.52846 & actuator bandwidth ($\Delta_3$) \\
0.0888 & 0.51625 & actuator bandwidth ($\Delta_3$) \\
0.0900 & none    & NMP ceiling ($\Delta_4$): unstable for all $\eta$ \\
\bottomrule
\end{tabular}
\end{table}

\begin{table}[h]\centering
\caption{Executed consistency checks. Agreement establishes implementation correctness only.}
\begin{tabular}{ll}
\toprule
Check & Result \\
\midrule
Symbolic $\det(sI-A)$ versus \eqref{eq:poly} & exact equality \\
Normalised nonlinear Jacobian versus \eqref{eq:A} & max error $1.67\times10^{-13}$ \\
Prop.~\ref{prop:att} against sampled ratios & matches Routh condition \\
Prop.~\ref{prop:five}: abscissa with $\chi\in\{0.4,1,3.3\}$, $g_x,g_v$ fixed & identical to $10^{-12}$ \\
Prop.~\ref{prop:ceiling} at three parameter sets, $\eta=400$ & sign change across $g_x^{\max}$ \\
Matched-group normalised trajectories & max difference $1.24\times10^{-12}$ \\
Exact static mixing, synthetic equal-axis case & pole-abscissa spread $<1.5\times10^{-15}$ \\
\bottomrule
\end{tabular}
\end{table}

\section{Limitations}

HIL validates loaded actuator influence on an emulated cascade, not physical rocket translation; a
fixed pivot, a suspension, HIL, and free flight obey different equations and support different
claims. The reduced model excludes flexible structure, gimbal moving-mass reaction, estimator
dynamics, and environmental disturbance until identified; these may dominate at small scale. An IMU
measures specific force, so acceleration is not a gravity reference during thrust --- encoder truth is
mandatory. A stability boundary derived for one controller is not a controller-independent property
of the vehicle. Equation~\eqref{eq:law} is one candidate form among several and may be rejected. The
prior-art audit is public-source and non-exhaustive, and the patent observation is not a
freedom-to-operate opinion.

\section{Conclusion}

The analytical baseline is now exact and tighter than the small-scale TVC literature it draws on: a
five-group reduction in which inertia ratio is not independent, a closed-form non-minimum-phase gain
ceiling roughly an order of magnitude below the usual necessary condition, and an explicit
two-mechanism description of the stable region. The experimental contribution is deliberately narrow
and falsifiable: measure whether unloaded servo calibration transfers to a loaded two-axis gimbal,
and whether a one-parameter reflected-inertia correction restores held-out stability-boundary
prediction. The deliverables are an open dataset, a reproducible identification-and-prediction
pipeline, and a quantified transfer envelope for the servo class that essentially every small TVC
testbed uses.

\begin{thebibliography}{99}
\bibitem{spannagl2021} L.~Spannagl et al. Design, optimal guidance and control of a low-cost re-usable electric model rocket. \emph{IROS}, 6344--6351, 2021. arXiv:2103.04709.
\bibitem{linsen2022} R.~Linsen, P.~Listov, A.~de Lajarte, R.~Schwan, C.~N.~Jones. Optimal thrust vector control of an electric small-scale rocket prototype. \emph{ICRA}, 1996--2002, 2022.
\bibitem{erocket2026} P.~Santos, A.~Fonte, P.~Martins, P.~Oliveira. The E-Rocket: low-cost testbed for TVC rocket GNC validation, 2026. arXiv:2512.06535.
\bibitem{jacqmein1969} W.~A.~Jacqmein, D.~L.~Hall, M.~K.~Craig. A design parameter synthesis derived from a mathematical analysis of a hypothetical lunar flying vehicle. NASA MSC GWP 10,084, 1969. NTRS 19700025079.
\bibitem{correa2013} A.~E.~A.~Corr\^{e}a. Projeto de controlador proporcional derivativo para o terceiro est\'{a}gio de um ve\'{i}culo lan\c{c}ador. MSc thesis, Universidade de Bras\'{i}lia, 2013.
\bibitem{santos2024} P.~dos Santos, P.~Oliveira. Thrust vector control and state estimation architecture for low-cost small-scale launchers. \emph{Acta Astronautica} 222:52--68, 2024.
\bibitem{hauser1992} J.~Hauser, S.~Sastry, G.~Meyer. Nonlinear control design for slightly non-minimum phase systems: application to V/STOL aircraft. \emph{Automatica} 28(4):665--679, 1992.
\bibitem{orr2023} J.~S.~Orr et al. Advanced modeling of control-structure interaction in thrust vector control systems. AAS 23-153, 2023. NTRS 20230000427.
\bibitem{singh2025} E.~Singh. Design and testing of a low-cost 3D-printed servo gimbal for thrust vector control in model rockets, 2025. arXiv:2509.00061.
\bibitem{fonte2025} A.~Fonte, P.~Santos, P.~Oliveira. Control and navigation of a 2-D electric rocket, 2025. arXiv:2509.19970.
\bibitem{quadrocket2026} P.~Santos, J.~Reis, P.~Oliveira, C.~Silvestre. QuadRocket: an aerial robotic testbed for adaptive thrust-vector control of rocket-like vehicles, 2026. arXiv:2607.02474.
\bibitem{pidnet2026} C.~Calme, L.~Salley, A.~J.~Mu\~{n}oz-V\'{a}zquez, L.~F.~Zapata-Rivera. Robust PIDNet control of a dual-actuator thrust vectoring platform, 2026. arXiv:2607.27793.
\bibitem{ddobs2026} C.~Calme, L.~Salley, L.~F.~Zapata-Rivera, A.~J.~Mu\~{n}oz-V\'{a}zquez. Data-driven disturbance-observer-based actuator-space control of a dual-axis thrust-vectoring platform. \emph{Applied Sciences} 16(16):8330, 2026.
\bibitem{ep3529683} Thrust vectored multicopters. Patent EP3529683B1, 2021.
\bibitem{wo2025122161} Apparatus, system, and method to control torque or lateral thrust applied to a load suspended on a suspension cable. Patent WO2025122161A1, 2025.
\end{thebibliography}

\appendix
\section{Changes from revision 1}
\begin{enumerate}
\item The contribution is now a single falsifiable experimental hypothesis (H1, Eq.~\eqref{eq:law})
rather than a general ``validity envelope'' study. A generic dimensionless map is abandoned.
\item Proposition~\ref{prop:five} (five-group reduction) is promoted from a defensive remark to a
theorem, with its consequence for experiment design stated explicitly.
\item Proposition~\ref{prop:ceiling} replaces the loose necessary condition $\chi k_x<1$, which is
shown to be $\approx 11\times$ conservative, with an exact ceiling and its NMP interpretation.
\item The two-mechanism structure of the stable region is identified and tabulated.
\item Reaction loading is promoted from optional to mandatory; the unloaded-HIL default is rejected
because it makes H2 tautological.
\item The scan variable is changed from $\eta$ to a software gain, since servo bandwidth is not
continuously tunable on hardware.
\item A boundary estimator is specified (abscissa regression with bootstrap CI) where revision 1 had
none, and the primary outcome is changed from a binary label to abscissa prediction error.
\item Cross-axis coupling is given a predicted mechanism ($\sigma$, gyroscopic) instead of an
exploratory test.
\item Discrete time is made primary rather than a caveat.
\item Configuration count raised from 3 to $\geq 6$ per axis, with the unit of analysis stated.
\item Typographical: $\dot x$ restored in the physical-variable form of Eq.~\eqref{eq:cascade}.
\end{enumerate}

\end{document}
